%%%%%%%%%%%%%%%%%%%%%%%%%%%%%%%%%%%%%%%%%%%%%%%%%%%%%%%%%%%%%%%%%%%%%%%%%%%%%%%%
% conclusion.tex:
%%%%%%%%%%%%%%%%%%%%%%%%%%%%%%%%%%%%%%%%%%%%%%%%%%%%%%%%%%%%%%%%%%%%%%%%%%%%%%%%
\chapter{Conclusion and Discussion}
\label{conclusion_chapter}
%%%%%%%%%%%%%%%%%%%%%%%%%%%%%%%%%%%%%%%%%%%%%%%%%%%%%%%%%%%%%%%%%%%%%%%%%%%%%%%%

The work presented in this dissertation introduces two machine learning algorithms, \aframe and \amplfi, and documents their development, validation, and real-time deployment as an end-to-end gravitational wave search and parameter estimation pipeline.
The central demonstration is that machine learning methods can achieve sensitivity competitive with traditional matched-filter searches for compact binary coalescences while delivering alerts and skymaps at latencies that are orders of magnitude lower.
Together, \aframe and \amplfi represent the first demonstration of an end-to-end machine learning pipeline for gravitational waves operating in a live observing run.

Beyond these immediate results, this dissertation addresses a question that has remained open since the first applications of machine learning to gravitational-wave detection: whether such methods are mature enough to be trusted in a production environment alongside the established pipelines that have shaped the field.
Meeting this standard requires more than demonstrating competitive sensitivity in controlled experiments.
It requires handling the practical demands of continuous operation: data quality variation, non-stationary noise, fault tolerance, and multi-messenger data product generation.
By designing \aframe and \amplfi around these operational requirements and validating them against matched-filter searches on real observing data, this work provides a template for how machine learning pipelines can be developed and deployed responsibly in gravitational-wave astronomy.

\section{Summary of Contributions}

Chapter~\ref{chap:aframe_methods} introduced \aframe, a convolutional neural network trained to detect compact binary coalescence signals in two-detector strain data from LIGO.
Rather than evaluating a large template bank as in matched-filter searches, \aframe processes overlapping 1.5-second windows of coincident Hanford and Livingston strain data in real time to produce a scalar detection statistic.
The key technical contributions include a data augmentation strategy that enables the network to generalize to a wide range of signal morphologies and noise conditions, and an efficient GPU-based inference deployment capable of analyzing data continuously.
To enable direct comparison with matched-filter searches, we developed a sensitive volume metric that places ML and traditional pipelines on equal footing.
Evaluated on a test set drawn from the third observing run, \aframe identified 8 of 9 GWTC-3 astrophysical candidates at false alarm rates below 1 per year --- the minimum achievable given the test set duration --- while reporting zero non-catalog events at false alarm rates below 5 per month.

Chapter~\ref{chap:o3_catalog} presented the results of applying \aframe to the full O3 data set as an end-to-end search.
This analysis recovered 38 events from GWTC-3 with $p_\mathrm{astro} > 0.5$, together with 2 additional candidates from the IAS catalog and 1 from the OGC catalog --- a count comparable to that reported by individual matched-filter pipelines operating on the same data.
Of the 15 GWTC-3 candidates not recovered by \aframe, 5 had secondary masses below $5\,M_\odot$ and fell outside the training prior; the remaining 10 had network matched-filter SNRs below 12.6.
Parameter estimation was also performed on the IAS and OGC candidates using Bilby with the IMRPhenomXPHM waveform approximant, yielding posteriors broadly consistent with those reported by the respective catalogs.
The O3 search demonstrated competitive sensitivity to high-mass binary black hole mergers while also revealing the current model's limitations: reduced sensitivity at lower component masses, and a requirement for simultaneous operation of both LIGO detectors.

Chapter~\ref{chap:amplfi_methods} introduced \amplfi, a likelihood-free inference algorithm for rapid parameter estimation.
\amplfi uses a normalizing flow trained to estimate posterior distributions over source parameters --- including chirp mass, mass ratio, and sky location --- conditioned on strain data from both LIGO detectors.
The algorithm produces posterior samples in under two seconds from the time of a detection trigger, achieving an end-to-end latency of approximately six seconds including detection, significance evaluation, parameter estimation, and sky map generation.
A major technical contribution is the implementation of compact binary waveform generation directly on GPU, enabling on-the-fly waveform simulation during training at negligible computational cost.
\amplfi was designed for co-deployment with \aframe on shared GPU hardware, minimizing communication overhead and enabling the full pipeline latency demonstrated here.

Chapter~\ref{chap:amplfi_results} validated \amplfi against established methods using data from O3.
Comparisons with BAYESTAR, the standard low-latency sky localization algorithm, show that \amplfi produces sky maps of comparable quality for well-localized events.
Comparisons of recovered source parameter posteriors with those from GWTC-3 demonstrate broad consistency, including cases where posterior support extends near or beyond the boundaries of the training prior.

Chapter~\ref{chap:online} described the integration of \aframe and \amplfi into a production pipeline and the deployment of this pipeline during the final segment of LIGO's fourth observing run, O4c, from August 28 to November 25, 2025.
Of the 24 compact binary events for which data from both Hanford and Livingston were available, \aframe detected 23 (96\%).
The single non-detection was due to a computing infrastructure failure that prevented data delivery to the pipeline, rather than a limitation of the algorithm.
For 9 events where \aframe was the preferred pipeline at the time of the preliminary alert, \amplfi sky maps were distributed to the astronomical community via the General Coordinates Network.
The deployment relied on purpose-built software infrastructure, including the \mlgw library (Appendix~\ref{app:ml4gw}) for GPU-accelerated gravitational-wave data processing and a multiprocessing architecture that generates sky maps, source classifications, and significance estimates in parallel.
This deployment demonstrates that machine learning methods are mature enough for production use in gravitational wave astronomy and can contribute alongside traditional matched-filter pipelines in an operational setting.

\section{Outlook}

Several extensions of this work are already underway; others represent longer-horizon priorities.

On the detection side, the most significant limitation of the current \aframe implementation is the requirement that both LIGO Hanford and Livingston be operating simultaneously.
Development of multi-detector configurations --- supporting arbitrary detector subsets including three-detector, single-detector, and heterogeneous network modes --- will increase the fraction of observing time for which \aframe can contribute.
An extension of the training mass distribution to lower-mass systems, including binary neutron star and neutron star--black hole mergers, will make the search relevant for the events most likely to have electromagnetic counterparts and yield the most multi-messenger science.
Extending \aframe to BNS systems poses a fundamental signal-duration challenge: a 1.4+1.4\,$M_\odot$ BNS inspiral spends 60--120 seconds in the LIGO sensitive band, far exceeding the 1.5-second input windows on which the current model operates.
One experimental approach under investigation evaluates each input window in the frequency domain against a log-normally spaced grid of chirp mass hypotheses --- approximately 100 templates per window --- using heterodyning to shift each template to baseband, where the slowly evolving signal envelope can be matched efficiently with a compact inner product.
Whether this approach can be made computationally cheap enough for real-time operation remains an open question, but initial experiments are encouraging.

On the parameter estimation side, extending \amplfi to binary neutron star and neutron star--black hole systems is a natural next step given the scientific priority of multi-messenger follow-up.
The longer signal durations and richer parameter spaces of these systems will require dedicated training datasets and may motivate architectural changes to the normalizing flow.
A heterodyning-based approach analogous to that described above for detection is also being explored on the parameter estimation side, compressing long BNS signal representations into a form compatible with \amplfi's existing input window; this work is at an earlier stage than the detection-side development.
Future versions may also incorporate spin-precessing waveform models beyond IMRPhenomD and could support burst signal morphologies, enabling collaboration with unmodeled event search pipelines.

On the deployment side, integrating a distributed task orchestration framework such as Ray or Celery will improve fault tolerance by allowing failed sub-tasks to be retried independently without restarting the entire pipeline.
Several latency reductions have already been made following O4c.
The post-resampling crop was tightened from a conservative 1-second buffer to one based on the filter's transient length (typically less than 0.2 seconds), recovering resampling latency without modifying the underlying filter.
Also in active development is an auxiliary regression head for \aframe to directly output a merger time estimate, which would eliminate the top-hat integration step used to locate the coalescence time.
Longer-horizon improvements include replacing the current symmetric (zero-phase) whitening filter with a causal (minimum-phase) implementation to eliminate the latency introduced by the whitening step, and extending the detector network to include Virgo, which would broaden sky localization coverage for three-detector events.
Infrastructure improvements outside the pipeline will also reduce latency: the next-generation data delivery (NGDD) system, currently in development for future observing runs, will deliver data in 1/16-second chunks rather than 1-second frames, reducing the discrete packet latency, and will deliver each chunk immediately rather than waiting for all auxiliary channels to be ready, reducing the data preparation and transfer latency.

Looking further ahead, the next generation of ground-based detectors --- Cosmic Explorer and the Einstein Telescope --- are expected to detect hundreds of thousands of compact binary mergers per year at far greater distances than current instruments can reach.
At those event rates, machine learning methods will likely be essential infrastructure rather than merely an alternative to matched filtering.
The algorithms, software frameworks, and operational experience developed in this dissertation provide a foundation for meeting that challenge.

%%%%%%%%%%%%%%%%%%%%%%%%%%%%%%%%%%%%%%%%%%%%%%%%%%%%%%%%%%%%%%%%%%%%%%%%%%%%%%%%
